\section{Conclusion}
\label{sec:conclusion}

Cloud-native operations have made change cheap, fast, and voluminous---and have
thereby made the \emph{record} of change the system's most perishable component.
This paper named the resulting deferred cost evidence debt and did the work the
name requires: non-circular definitions grounded in observables; a formal model in
which debt is expected excess reconstruction cost, reconstruction is linkage
inference, and decay is a survival process; analytic results locating exactly where
the financial metaphor breaks (stochastic step interest, mandatory write-offs) and
what remains when it is removed; a metric suite honest enough to reject one of its
own commissioned members; an executed, fully reproducible degradation experiment
showing that redundancy absorbs small deficits, that deficit classes collapse
reconstruction superadditively, that effort inflates before failure appears, and
that reconstruction under heavy deficit fabricates confident false history at rates
that should alarm anyone who has trusted an after-the-fact timeline; a bounded,
age-stratified Argo CD observation showing that missing external issue linkage
occurs in public repository evidence while review-record and release-lineage
presence controls remain high; and a field protocol with stated refutation
conditions. The repository observation is not a measurement of intent quality,
decay, evidence debt, or reconstruction cost.

The central claim---deferral without elimination---emerges from this treatment
with its proper status: not an axiom but a structured hypothesis, true in the
simulation's closed world in a specific, mechanistic way (payment in effort, then
error, then loss), and awaiting field magnitudes. The constructs are offered for
reuse: a deficit census any organization can run, a decay schedule any organization
can read off its own retention settings, and a definition of debt that makes the
implicit question---\emph{what will it cost us to not know what we did?}---explicit,
estimable, and budgetable. Systems will keep forgetting; the choice organizations
retain is whether the forgetting is priced.
